% =====================================================================
%  ENTERPRISE USER MANUAL: CMS SEKOLAH
%  Dibuat oleh: Technical Writer AI
%  Versi: 1.0 | Tanggal: \today
% =====================================================================

\documentclass[11pt,a4paper,bahasa]{report}

% --- ENCODING & BAHASA ---
\usepackage[utf8]{inputenc}
\usepackage[T1]{fontenc}
\usepackage[bahasa]{babel}

% --- FONT ---
\usepackage{lmodern}

% --- LAYOUT & GEOMETRI ---
\usepackage[
  a4paper,
  top=2.5cm, bottom=2.5cm,
  left=3cm,  right=2.5cm,
  headheight=14pt
]{geometry}

% --- GAMBAR ---
\usepackage{graphicx}
\usepackage{float}

% --- WARNA ---
\usepackage[dvipsnames,table]{xcolor}
\definecolor{primaryblue}{RGB}{14,165,233}
\definecolor{slategray}{RGB}{71,85,105}
\definecolor{lightbg}{RGB}{248,250,252}
\definecolor{warnyellow}{RGB}{250,176,5}

% --- TABEL ---
\usepackage{booktabs}
\usepackage{tabularx}
\usepackage{array}
\usepackage{longtable}

% --- HYPERLINK ---
\usepackage[
  colorlinks=true,
  linkcolor=slategray,
  urlcolor=primaryblue,
  citecolor=slategray,
  bookmarks=true,
  bookmarksnumbered=true
]{hyperref}

% --- HEADER / FOOTER ---
\usepackage{fancyhdr}
\pagestyle{fancy}
\fancyhf{}
\fancyhead[L]{\small\textcolor{slategray}{\nouppercase{\leftmark}}}
\fancyhead[R]{\small\textcolor{slategray}{CMS Sekolah -- Manual Pengguna}}
\fancyfoot[C]{\small\thepage}
\renewcommand{\headrulewidth}{0.4pt}

% --- JUDUL BAB ---
\usepackage{titlesec}
\titleformat{\chapter}[display]
  {\normalfont\huge\bfseries\color{slategray}}
  {\chaptertitlename\ \thechapter}{12pt}{\Huge}

% --- DAFTAR (LIST) ---
\usepackage{enumitem}
\setlist[enumerate]{leftmargin=*, itemsep=2pt}
\setlist[itemize]{leftmargin=*, itemsep=2pt}

% --- KOTAK PERINGATAN / CATATAN ---
\usepackage{tcolorbox}
\tcbuselibrary{skins,breakable}

\newtcolorbox{notebox}[1][Catatan]{
  colback=lightbg, colframe=primaryblue,
  fonttitle=\bfseries, title=#1,
  breakable, enhanced
}
\newtcolorbox{warnbox}[1][Perhatian]{
  colback=yellow!10, colframe=warnyellow,
  fonttitle=\bfseries, title=#1,
  breakable, enhanced
}

% --- KODE ---
\usepackage{listings}
\lstset{
  basicstyle=\ttfamily\small,
  backgroundcolor=\color{lightbg},
  frame=single,
  breaklines=true
}

% --- URL PANJANG ---
\usepackage{url}
\urlstyle{same}
\usepackage{microtype}

% --- SPACING ---
\usepackage{parskip}
\setlength{\parindent}{0pt}

% --- METADATA DOKUMEN ---
\hypersetup{
  pdftitle={Manual Pengguna CMS Sekolah},
  pdfauthor={Tim Teknologi Sekolah},
  pdfsubject={Panduan Administrasi dan Penggunaan CMS Sekolah},
  pdfkeywords={CMS, Sekolah, Manual, Laravel, Admin}
}

% =====================================================================
\begin{document}
\sloppy

% --- HALAMAN JUDUL ---
\begin{titlepage}
  \centering
  \vspace*{2cm}

  {\Huge\bfseries\color{slategray} Manual Pengguna Sistem\\[0.4em]
  \textcolor{primaryblue}{CMS Sekolah}}

  \vspace{1cm}
  {\large Panduan Lengkap Administrasi Backend}

  \vspace{0.5cm}
  \rule{0.6\textwidth}{0.4pt}

  \vspace{1.5cm}
  \begin{tabular}{rl}
    \textbf{Versi}    & 1.0 \\
    \textbf{Tanggal}  & \today \\
    \textbf{Platform} & Laravel 11 + Livewire \\
    \textbf{Akses}    & \texttt{/admin/dashboard} \\
  \end{tabular}

  \vfill
  {\small Dokumen ini bersifat \textit{internal} dan ditujukan bagi Administrator,
  Editor, dan pengelola website sekolah.}
\end{titlepage}

% --- DAFTAR ISI ---
\tableofcontents
\clearpage

% --- BAB-BAB ---
% =====================================================================
%  BAB 1: PENDAHULUAN & AKSES SISTEM
% =====================================================================

\chapter{Pendahuluan \& Akses Sistem}

CMS Sekolah adalah sistem manajemen konten (\textit{Content Management System}) berbasis web
yang dibangun di atas platform \textbf{Laravel 11} dan \textbf{Livewire}. Sistem ini dirancang
khusus untuk memenuhi kebutuhan institusi pendidikan dalam mengelola informasi publik sekolah
secara mandiri, aman, dan efisien---tanpa membutuhkan keahlian pemrograman.

Sistem ini mencakup modul-modul seperti Berita, Pengumuman, Agenda, Galeri, Data Kepegawaian,
Manajemen Media, hingga Konfigurasi Website. Seluruh perubahan yang dilakukan melalui panel
admin akan langsung terrefleksi pada halaman publik website sekolah secara \textit{real-time}.

% -------------------------------------------------------------------
\section{Matriks Peran (Role) Pengguna}
% -------------------------------------------------------------------

Akses ke panel admin dibatasi oleh tiga tingkatan peran yang berbeda:

\begin{table}[H]
\centering
\begin{tabularx}{\textwidth}{|l|X|}
\toprule
\textbf{Peran (Role)} & \textbf{Hak Akses} \\
\midrule
\textbf{Super Admin} &
  Akses penuh ke \textbf{seluruh} modul termasuk: Manajemen Pengguna,
  Pengaturan Global, URL Login Kustom, Pods (tipe konten dinamis),
  SEO, dan konfigurasi tampilan. \\
\midrule
\textbf{Admin} &
  Dapat mengelola semua konten (Berita, Pengumuman, Agenda, Galeri,
  Slider, Media, Unduhan, Staf, Fasilitas, Prestasi, Ekstrakurikuler).
  Tidak bisa mengubah Pengaturan inti dan tidak bisa mengelola akun pengguna lain. \\
\midrule
\textbf{Editor} &
  Hanya dapat membuat, mengedit, dan menerbitkan Berita (\textit{Post}).
  Tidak memiliki akses ke modul lain. \\
\bottomrule
\end{tabularx}
\caption{Matriks hak akses pengguna CMS Sekolah}
\label{table:roles}
\end{table}

% -------------------------------------------------------------------
\section{Skenario \& Alur Kerja: Login Pertama Kali}
% -------------------------------------------------------------------

Misalkan Ibu Dewi baru saja ditunjuk sebagai Admin website oleh Kepala Sekolah.
\textbf{Super Admin} (Bapak IT Sekolah) telah membuatkan akun untuknya dengan
\textit{role Admin}.

\begin{enumerate}
  \item Bapak IT memberi tahu Ibu Dewi alamat URL login khusus sekolah, misalnya:
        \texttt{http://situs-sekolah.sch.id/\textbf{portal-admin}}.
  \item Ibu Dewi membuka URL tersebut di peramban dan muncul halaman login.
  \item Ia memasukkan \textbf{Email} dan \textbf{Password} sementara yang diberikan.
  \item Sistem memverifikasi kredensial. Jika salah 5 kali berturut-turut,
        sistem secara otomatis memblokir percobaan login selama 1 menit
        (\textit{rate limiting: 5 percobaan per menit}).
  \item Jika berhasil, Ibu Dewi diarahkan ke halaman \textbf{Dashboard} (\texttt{/admin/dashboard}).
\end{enumerate}

\begin{figure}[H]
  \centering
  \includegraphics[width=0.9\textwidth]{../Img/Backend/screencapture-127-0-0-1-8000-admin-dashboard-2026-06-16-17_51_15.png}
  \caption{Halaman Dashboard Admin --- Ringkasan Statistik Konten}
  \label{fig:dashboard}
\end{figure}

% -------------------------------------------------------------------
\section{Panduan Penggunaan: Login \& Dashboard}
% -------------------------------------------------------------------

\subsection{Proses Login}

\begin{enumerate}
  \item Buka peramban web dan navigasikan ke URL Login yang diberikan oleh Super Admin
        (defaultnya: \texttt{/login}, namun bisa dikustomisasi).
  \item Isi kolom \textbf{Email} dengan alamat surel akun Anda.
  \item Isi kolom \textbf{Password} dengan kata sandi.
  \item Tekan tombol \textbf{Masuk}.
\end{enumerate}

\begin{notebox}[Informasi Keamanan]
  Jika URL Login sudah diubah oleh Super Admin (misalnya menjadi \texttt{/portal-guru}),
  akses melalui \texttt{/login} standar akan menghasilkan halaman \textit{404 Not Found}.
  Pastikan Anda menyimpan tautan login yang benar di bookmark peramban Anda.
\end{notebox}

\subsection{Memahami Dashboard}

Setelah login, Anda akan melihat halaman \textbf{Ringkasan Admin} yang terdiri dari:

\begin{itemize}
  \item \textbf{Kartu Statistik (8 buah)}: Menampilkan total data aktif dari modul:
        \begin{itemize}
          \item \textbf{Berita} --- total artikel yang tersimpan (semua status)
          \item \textbf{Pengumuman} --- total pengumuman
          \item \textbf{Agenda} --- total agenda/kegiatan
          \item \textbf{Galeri} --- total album galeri
          \item \textbf{Program Unggulan} --- total ekstrakurikuler
          \item \textbf{Prestasi} --- total prestasi siswa
          \item \textbf{Staff/Guru} --- total data kepegawaian
          \item \textbf{Media (Gambar)} --- total file gambar di pustaka media
        \end{itemize}
  \item \textbf{Panel Berita Terbaru}: Daftar 5 berita terakhir. Setiap entri menampilkan
        judul, status (\textit{Draft / Published}), dan waktu pembaruan terakhir.
        Terdapat tautan \textbf{Edit} untuk langsung membuka form pengeditan.
  \item \textbf{Panel Pengumuman Terbaru}: Mirip dengan Berita Terbaru, menampilkan
        pengumuman terakhir beserta statusnya.
  \item \textbf{Panel Agenda Mendatang}: Menampilkan agenda yang akan datang,
        diurutkan berdasarkan tanggal \textit{start\_at}.
  \item \textbf{Panel Pengguna Baru}: Menampilkan akun pengguna yang baru ditambahkan
        beserta alamat email mereka (hanya terlihat oleh Super Admin/Admin).
\end{itemize}

% -------------------------------------------------------------------
\section{Logout / Keluar dari Sistem}
% -------------------------------------------------------------------

Untuk keluar dari sesi admin:
\begin{enumerate}
  \item Klik nama pengguna atau ikon profil di sudut kanan atas \textit{sidebar}.
  \item Klik tombol \textbf{Keluar} / \textbf{Logout}.
  \item Sistem akan mengakhiri sesi dan mengarahkan Anda kembali ke halaman Login.
\end{enumerate}

\begin{warnbox}[Keamanan Sesi]
  Jangan pernah meninggalkan sesi admin terbuka di komputer umum atau bersama.
  Selalu \textbf{Logout} setelah selesai bekerja.
\end{warnbox}

% -------------------------------------------------------------------
\section{Penanganan Masalah Umum (Akses \& Login)}
% -------------------------------------------------------------------

\begin{itemize}
  \item \textbf{Halaman Login Menampilkan 404}: Super Admin telah mengubah URL Login
        dari \texttt{/login} menjadi alamat rahasia lain. Hubungi Super Admin untuk
        mendapatkan tautan yang benar.
  \item \textbf{Akun Diblokir / Tidak Bisa Login}: Sistem menerapkan pembatasan
        5 percobaan login per menit. Tunggu 1 menit lalu coba kembali. Jika lupa
        password, hubungi Super Admin untuk melakukan reset.
  \item \textbf{Halaman Dashboard Kosong / Angka Statistik Nol}: Ini normal untuk
        sistem yang baru dipasang. Mulailah mengisi konten melalui modul-modul
        yang tersedia.
\end{itemize}

% =====================================================================
%  BAB 2: MANAJEMEN KONTEN UTAMA
%  (Berita / Posts, Pengumuman / Announcements, Agenda / Events)
% =====================================================================

\chapter{Manajemen Konten Utama}

Modul ini mencakup tiga jenis konten informasi publik paling penting
dalam website sekolah: \textbf{Berita}, \textbf{Pengumuman}, dan \textbf{Agenda}.
Ketiganya dapat dikelola oleh pengguna berperan \textbf{Admin} atau \textbf{Editor}
(khusus Berita).

% -------------------------------------------------------------------
\section{Modul Berita (\textit{Posts})}
% -------------------------------------------------------------------

Modul Berita adalah sistem penerbitan artikel berbasis \textit{WYSIWYG editor}
(menggunakan \textbf{Quill Editor}). Setiap artikel mendukung fitur kategori,
\textit{tag}, thumbnail, SEO, penjadwalan tayang otomatis, dan penandaan sebagai
artikel unggulan.

\begin{figure}[H]
  \centering
  \includegraphics[width=0.9\textwidth]{../Img/Backend/screencapture-127-0-0-1-8000-admin-posts-2026-06-16-17_51_29.png}
  \caption{Daftar Berita di Panel Admin}
  \label{fig:posts-index}
\end{figure}

\subsection{Skenario: Menerbitkan Berita Prestasi Sekolah}

Misalkan SMA Nusantara baru saja meraih Juara 1 Olimpiade Sains Nasional.
Guru Pembimbing (berperan \textit{Editor}) menulis draf, kemudian Admin
menambahkan thumbnail dan memublikasikannya.

\begin{figure}[H]
  \centering
  \includegraphics[width=0.9\textwidth]{../Img/Backend/screencapture-127-0-0-1-8000-admin-posts-22-edit-2026-06-16-17_51_38.png}
  \caption{Form Editor Berita --- Tampilan Lengkap}
  \label{fig:posts-edit}
\end{figure}

\subsection{Panduan: Menulis Berita Baru}

\begin{enumerate}
  \item Klik menu \textbf{Berita} di \textit{sidebar} kiri, lalu klik tombol \textbf{Tulis Berita Baru}.
  \item Pada panel kiri atas, isi kolom \textbf{Judul} (contoh: ``Juara 1 OSN Bidang Matematika'').
        Kolom \textbf{Slug} akan terisi otomatis berdasarkan judul (dapat diubah manual jika diperlukan).
  \item Tuliskan isi lengkap artikel pada area \textbf{Konten} menggunakan Quill Editor.
        Editor mendukung pemformatan teks tebal, miring, daftar, tautan, dan penyisipan gambar.
  \item Isi kolom \textbf{Excerpt} --- ringkasan singkat artikel (1-3 kalimat) yang
        akan tampil di halaman daftar berita.
  \item Pada panel kanan, atur \textbf{Status} melalui \textit{dropdown}:
        \begin{itemize}
          \item \textit{Draf} --- tersimpan, tidak terlihat publik
          \item \textit{Review} --- menunggu persetujuan Admin
          \item \textit{Terjadwal} --- tayang otomatis pada waktu yang ditentukan
          \item \textit{Terbit} --- langsung terlihat publik
          \item \textit{Arsip} --- disembunyikan dari publik tanpa dihapus
        \end{itemize}
  \item Jika memilih \textit{Terjadwal}, isi \textbf{Jadwal Tayang} (format: Tanggal dan Jam).
  \item Centang \textbf{Tandai sebagai unggulan} jika artikel ini harus tampil di
        \textit{carousel} utama halaman depan.
  \item Pilih \textbf{Kategori} dari \textit{dropdown} yang tersedia.
  \item Tambahkan \textbf{Tag}: ketikkan nama \textit{tag} pada kolom input lalu klik tombol \textbf{Tambah},
        atau klik salah satu \textit{tag} yang sudah ada untuk memilihnya. Tag yang dipilih akan muncul
        sebagai \textit{badge} biru yang bisa dihapus dengan klik tanda \textbf{×}.
  \item Pilih gambar \textbf{Thumbnail}: klik tombol \textbf{Pilih Media} untuk membuka
        Pustaka Media. Disarankan menggunakan rasio \textbf{16:9} untuk hasil terbaik.
        Klik \textbf{Hapus} untuk membatalkan pilihan thumbnail.
  \item (Opsional) Pilih \textbf{OG Image} (Open Graph Image) untuk pratinjau saat artikel
        dibagikan di media sosial. Disarankan rasio \textbf{1.91:1}.
  \item Isi panel \textbf{SEO} di bagian bawah: \textit{Meta Title}, \textit{Meta Description},
        \textit{Canonical URL}. Aktifkan \textit{No Index} jika tidak ingin artikel ini diindeks
        oleh mesin pencari.
  \item Klik salah satu tombol aksi di \textit{action bar} yang menempel di bagian bawah layar:
        \begin{itemize}
          \item \textbf{Simpan Draf} (hitam) --- menyimpan sebagai draf
          \item \textbf{Simpan Perubahan} (biru) --- menyimpan dengan status yang dipilih
          \item \textbf{Publikasikan} (biru border) --- langsung mengubah status menjadi \textit{Published}
          \item \textbf{Preview} --- membuka pratinjau artikel di tab baru
        \end{itemize}
\end{enumerate}

\begin{notebox}[Izinkan Komentar]
  Centang opsi \textbf{Izinkan komentar} jika Anda ingin mengaktifkan fitur
  komentar pembaca pada artikel ini. Fitur ini bersifat per-artikel.
\end{notebox}

\subsection{Mengelola Kategori \& Tag}

\textbf{Kategori} dan \textbf{Tag} dikelola melalui menu terpisah di \textit{sidebar}:

\begin{enumerate}
  \item Klik menu \textbf{Kategori} untuk menambah, mengedit, atau menghapus kategori.
        Setiap kategori memiliki nama dan \textit{slug} unik.
  \item Klik menu \textbf{Tag} untuk mengelola tag. Tag bisa juga ditambahkan
        langsung dari dalam form penulisan berita.
\end{enumerate}

\subsection{Penanganan Masalah: Berita}

\begin{itemize}
  \item \textbf{Artikel tidak muncul di halaman publik}: Periksa apakah
        \textbf{Status} sudah diatur ke \textit{Terbit}. Status \textit{Draf},
        \textit{Review}, dan \textit{Arsip} tidak terlihat oleh publik.
  \item \textbf{Gambar Thumbnail tidak berubah}: Coba \textit{hard refresh} halaman
        (Ctrl+Shift+R). Perubahan gambar memerlukan pembersihan \textit{cache} browser.
  \item \textbf{Tombol Publikasikan tidak berfungsi}: Pastikan kolom \textbf{Judul}
        dan \textbf{Konten} sudah diisi. Kedua kolom ini bersifat wajib (\textit{required}).
\end{itemize}

% -------------------------------------------------------------------
\section{Modul Pengumuman (\textit{Announcements})}
% -------------------------------------------------------------------

Modul Pengumuman digunakan untuk menerbitkan pemberitahuan resmi sekolah,
seperti jadwal libur, informasi ujian, atau edaran penting. Pengumuman mendukung
fitur \textit{Pin} (sematkan) dan lampiran dokumen.

\begin{figure}[H]
  \centering
  \includegraphics[width=0.9\textwidth]{../Img/Backend/screencapture-127-0-0-1-8000-admin-announcements-2026-06-16-17_51_50.png}
  \caption{Daftar Pengumuman di Panel Admin}
  \label{fig:announcements-index}
\end{figure}

\begin{figure}[H]
  \centering
  \includegraphics[width=0.9\textwidth]{../Img/Backend/screencapture-127-0-0-1-8000-admin-announcements-1-edit-2026-06-16-17_52_05.png}
  \caption{Form Edit Pengumuman}
  \label{fig:announcements-edit}
\end{figure}

\subsection{Skenario: Menerbitkan Pengumuman Libur Nasional}

Kepala Sekolah meminta Admin TU untuk segera menerbitkan pengumuman bahwa sekolah
libur pada tanggal 17 Agustus. Admin membuat pengumuman baru, menyematkannya
(\textit{Pin}) agar muncul paling atas, dan melampirkan surat edaran PDF.

\subsection{Panduan: Membuat Pengumuman}

\begin{enumerate}
  \item Buka menu \textbf{Pengumuman} lalu klik tombol \textbf{Buat Pengumuman}.
  \item Isi kolom \textbf{Judul} (contoh: ``Libur HUT RI ke-80'').
  \item Pilih \textbf{Kategori} dari \textit{dropdown} (misal: ``Akademik'', ``Umum'').
        Pilih kosong jika tidak berkategori.
  \item Tentukan rentang aktif pengumuman:
        \begin{itemize}
          \item \textbf{Mulai} --- tanggal pengumuman mulai aktif (format: YYYY-MM-DD)
          \item \textbf{Selesai} --- tanggal pengumuman berakhir
        \end{itemize}
  \item Atur \textbf{Status}:
        \begin{itemize}
          \item \textit{Draft} --- belum terlihat publik
          \item \textit{Published} --- langsung terlihat publik
        \end{itemize}
  \item Centang \textbf{Pin Pengumuman} jika pengumuman ini harus muncul paling
        atas dalam daftar dan diberikan sorotan khusus.
  \item Tulis isi pengumuman di editor \textbf{Konten} (Quill Editor, tinggi 280px).
  \item (Opsional) Unggah berkas di kolom \textbf{Lampiran (PDF/DOC)}.
        Jika pengumuman sudah pernah memiliki lampiran dan Anda ingin menggantinya,
        centang kotak \textbf{Hapus lampiran} sebelum mengunggah file baru.
  \item Klik tombol hitam \textbf{Simpan Pengumuman}.
\end{enumerate}

\subsection{Penanganan Masalah: Pengumuman}

\begin{itemize}
  \item \textbf{Pengumuman hilang dari website sebelum tanggal selesai}: Pastikan
        \textbf{Status} masih \textit{Published}. Jika Anda atau orang lain tidak sengaja
        mengubah status ke \textit{Draft}, pengumuman akan tersembunyi.
  \item \textbf{File lampiran gagal diunggah}: Pastikan file berformat PDF atau DOC/DOCX
        dan ukurannya tidak melebihi batas unggah server (biasanya 2 MB untuk \textit{shared hosting}).
\end{itemize}

% -------------------------------------------------------------------
\section{Modul Agenda (\textit{Events})}
% -------------------------------------------------------------------

Modul Agenda mencatat jadwal kegiatan sekolah (seminar, lomba, ujian, perlombaan, dll.)
beserta lokasi, durasi, dan gambar acara.

\begin{figure}[H]
  \centering
  \includegraphics[width=0.9\textwidth]{../Img/Backend/screencapture-127-0-0-1-8000-admin-events-2026-06-16-17_52_13.png}
  \caption{Daftar Agenda di Panel Admin}
  \label{fig:events-index}
\end{figure}

\begin{figure}[H]
  \centering
  \includegraphics[width=0.9\textwidth]{../Img/Backend/screencapture-127-0-0-1-8000-admin-events-1-edit-2026-06-16-17_52_26.png}
  \caption{Form Edit Agenda --- Detail Waktu dan Lokasi}
  \label{fig:events-edit}
\end{figure}

\subsection{Panduan: Menambahkan Agenda}

\begin{enumerate}
  \item Masuk ke menu \textbf{Agenda} lalu klik tombol tambah.
  \item Isi \textbf{Judul} agenda (contoh: ``Ujian Tengah Semester Ganjil 2025'').
  \item Isi \textbf{Lokasi} tempat pelaksanaan (contoh: ``Aula SMA Nusantara'').
  \item Tentukan waktu mulai pada kolom \textbf{Mulai} (format: tanggal dan jam, contoh: 2025-10-14T08:00).
  \item Tentukan waktu selesai pada kolom \textbf{Selesai} (format sama).
  \item Tulis deskripsi lengkap agenda pada editor \textbf{Deskripsi}.
  \item Atur \textbf{Status}: \textit{Draft} atau \textit{Published}.
  \item Centang \textbf{Agenda Unggulan} jika kegiatan ini berskala besar dan
        ingin ditampilkan secara mencolok di halaman agenda publik.
  \item Unggah gambar/poster acara pada kolom \textbf{Gambar Agenda} melalui komponen
        \textit{image-upload}.
  \item Klik tombol hitam \textbf{Simpan}.
\end{enumerate}

\subsection{Penanganan Masalah: Agenda}

\begin{itemize}
  \item \textbf{Agenda tidak muncul di halaman depan}: Halaman depan biasanya menampilkan
        agenda yang memiliki tanggal \textit{Mulai} di masa mendatang. Agenda yang sudah
        lewat (\textit{expired}) tidak akan ditampilkan di \textit{upcoming events}.
  \item \textbf{Format waktu salah}: Kolom \textit{Mulai} dan \textit{Selesai} menggunakan
        format \texttt{YYYY-MM-DDTHH:mm}. Pastikan browser Anda menampilkan
        \textit{datetime-local picker} dengan benar.
\end{itemize}

% =====================================================================
%  BAB 3: MANAJEMEN KEPEGAWAIAN & AKADEMIK
%  (Staf/Guru, Prestasi, Ekstrakurikuler, Fasilitas)
% =====================================================================

\chapter{Manajemen Kepegawaian \& Data Akademik}

Modul ini menyediakan fitur untuk mengelola data sumber daya manusia sekolah
(guru, tenaga kependidikan, kepala sekolah) dan informasi akademik pendukung
(prestasi siswa, program ekstrakurikuler, fasilitas sekolah).

% -------------------------------------------------------------------
\section{Modul Staf \& Guru (\textit{Staffs})}
% -------------------------------------------------------------------

Data Staf dan Guru akan tampil secara dinamis di halaman publik website sekolah,
meliputi halaman \texttt{/guru}, \texttt{/tenaga-didik}, dan \texttt{/kepala-sekolah}.
Sistem membedakan tampilan berdasarkan kolom \textbf{Tipe} yang dipilih saat pendataan.

\begin{figure}[H]
  \centering
  \includegraphics[width=0.9\textwidth]{../Img/Backend/screencapture-127-0-0-1-8000-admin-staffs-2026-06-16-17_53_39.png}
  \caption{Daftar Staf \& Guru di Panel Admin}
  \label{fig:staffs-index}
\end{figure}

\begin{figure}[H]
  \centering
  \includegraphics[width=0.9\textwidth]{../Img/Backend/screencapture-127-0-0-1-8000-admin-staffs-1-edit-2026-06-16-17_53_52.png}
  \caption{Form Edit Data Guru --- Termasuk Sub-Panel Kepala Sekolah}
  \label{fig:staffs-edit}
\end{figure}

\subsection{Skenario: Mendaftarkan Guru Baru}

Pak Budi, S.Pd, baru saja resmi bergabung sebagai guru Matematika.
Admin kepegawaian perlu mendaftarkan profilnya agar muncul di halaman
daftar guru website sekolah.

\subsection{Panduan: Menambahkan Data Staf / Guru}

\begin{enumerate}
  \item Buka menu \textbf{Staf/Guru} di \textit{sidebar}, lalu klik tombol tambah.
  \item Pilih \textbf{Tipe} dari \textit{dropdown}:
        \begin{itemize}
          \item \textit{Guru} --- tampil di halaman \texttt{/guru}
          \item \textit{Tenaga Kependidikan} --- tampil di halaman \texttt{/tenaga-didik}
          \item \textit{Kepala Sekolah} --- tampil di halaman \texttt{/kepala-sekolah}
                dengan tampilan profil khusus yang lebih lengkap
        \end{itemize}
  \item Isi \textbf{Nama Lengkap} beserta gelar akademik.
  \item Isi \textbf{NIP / NIY} (Nomor Induk Pegawai / Yayasan).
  \item Pilih \textbf{Jabatan Struktur (Organisasi)} dari \textit{dropdown} jika pegawai ini
        memegang posisi struktural (misal: Wakil Kepala Sekolah Bidang Kurikulum).
        Pilih \textit{--- Tidak ada jabatan terstruktur ---} untuk guru biasa.
  \item Isi \textbf{Tugas / Jabatan Lainnya} secara bebas (misal: ``Wali Kelas XII IPA 1'').
  \item Isi \textbf{Mata Pelajaran Diampu} (misal: ``Matematika Wajib'').
  \item Isi \textbf{Pendidikan Terakhir} (misal: ``S1 Pendidikan Matematika, UNY'').
  \item Tentukan \textbf{Urutan (Tampil di Web)}: angka terkecil akan tampil paling atas.
        Gunakan angka 1 untuk yang paling prioritas.
  \item Tulis \textbf{Bio Singkat} di \textit{textarea} yang tersedia.

  \item Jika \textbf{Tipe} dipilih \textit{Kepala Sekolah}, isi juga sub-panel
        \textbf{Informasi Tambahan}:
        \begin{itemize}
          \item \textbf{Tempat Tanggal Lahir} (contoh: ``Palembang, 6 Juli 1975'')
          \item \textbf{Jabatan Fungsional} (contoh: ``Guru Ahli Madya (IV/b)'')
          \item \textbf{Status / Pangkat / GOL} (contoh: ``IV/b -- Pembina Tingkat I'')
          \item \textbf{Perangkat Daerah}
          \item \textbf{Unit Kerja} (contoh: ``SMA Negeri 1 Kepohbaru'')
        \end{itemize}

  \item Atur \textbf{Status Aktif}: pilih \textit{Aktif} untuk guru yang masih berdinas,
        \textit{Nonaktif} untuk yang sudah tidak aktif (data tetap tersimpan).
  \item Unggah \textbf{Foto} pas foto guru (komponen \textit{image-upload}).
  \item Klik tombol hitam \textbf{Simpan}.
\end{enumerate}

\begin{warnbox}[Penting: Tipe Kepala Sekolah]
  Halaman publik \texttt{/kepala-sekolah} dirancang khusus dan hanya akan menampilkan
  data staf yang kolom \textbf{Tipe}-nya diatur ke \textit{Kepala Sekolah}. Salah memilih
  tipe akan menyebabkan profil tidak muncul di halaman yang tepat.
\end{warnbox}

\subsection{Mengelola Jabatan Struktural (Posisi)}

Daftar jabatan yang tersedia di \textit{dropdown} \textbf{Jabatan Struktur} dikelola
melalui menu \textbf{Posisi (Jabatan)}:

\begin{enumerate}
  \item Klik menu \textbf{Posisi} di \textit{sidebar}.
  \item Tambahkan jabatan baru (misal: ``Wakil Kepala Sekolah Bidang Humas'').
  \item Jabatan yang baru ditambahkan akan langsung tersedia di \textit{dropdown}
        form staf.
\end{enumerate}

\subsection{Penanganan Masalah: Staf \& Guru}

\begin{itemize}
  \item \textbf{Urutan guru berantakan}: Pastikan setiap guru memiliki angka
        \textbf{Urutan} yang berbeda. Angka yang sama akan menghasilkan urutan acak.
  \item \textbf{Foto guru tidak muncul}: Pastikan proses unggah berhasil
        (ukuran file maksimal sesuai konfigurasi server). Format yang disarankan:
        JPG atau PNG, rasio 3:4 (potret).
\end{itemize}

% -------------------------------------------------------------------
\section{Modul Prestasi (\textit{Achievements})}
% -------------------------------------------------------------------

Modul Prestasi mendokumentasikan pencapaian siswa dalam berbagai kompetisi
atau kegiatan akademik/non-akademik.

\begin{figure}[H]
  \centering
  \includegraphics[width=0.9\textwidth]{../Img/Backend/screencapture-127-0-0-1-8000-admin-achievements-2026-06-16-17_54_46.png}
  \caption{Daftar Prestasi di Panel Admin}
  \label{fig:achievements-index}
\end{figure}

\subsection{Panduan: Menambahkan Data Prestasi}

\begin{enumerate}
  \item Buka menu \textbf{Prestasi} dan klik tambah.
  \item Isi nama judul prestasi, tingkatan (Kabupaten/Kota, Provinsi, Nasional,
        Internasional), tahun, dan nama siswa atau tim yang berprestasi.
  \item Unggah foto dokumentasi (piala, medali, sertifikat) sebagai gambar pendukung.
  \item Klik \textbf{Simpan}.
\end{enumerate}

% -------------------------------------------------------------------
\section{Modul Ekstrakurikuler (\textit{Extracurriculars})}
% -------------------------------------------------------------------

Menampilkan daftar program unggulan dan ekstrakurikuler sekolah di halaman
\texttt{/program-unggulan}.

\begin{figure}[H]
  \centering
  \includegraphics[width=0.9\textwidth]{../Img/Backend/screencapture-127-0-0-1-8000-admin-extracurriculars-2026-06-16-17_54_23.png}
  \caption{Daftar Ekstrakurikuler di Panel Admin}
  \label{fig:extracurriculars-index}
\end{figure}

\begin{figure}[H]
  \centering
  \includegraphics[width=0.9\textwidth]{../Img/Backend/screencapture-127-0-0-1-8000-admin-extracurriculars-1-edit-2026-06-16-17_54_38.png}
  \caption{Form Edit Ekstrakurikuler}
  \label{fig:extracurriculars-edit}
\end{figure}

\subsection{Panduan: Menambahkan Ekstrakurikuler}

\begin{enumerate}
  \item Buka menu \textbf{Program Unggulan} dan klik tambah.
  \item Isi nama ekstrakurikuler, deskripsi kegiatan (menggunakan editor teks),
        hari/waktu latihan, dan nama pembimbing.
  \item Unggah foto kegiatan sebagai gambar utama.
  \item Klik \textbf{Simpan}.
\end{enumerate}

% -------------------------------------------------------------------
\section{Modul Fasilitas (\textit{Facilities})}
% -------------------------------------------------------------------

Mendokumentasikan fasilitas fisik sekolah yang ditampilkan di halaman
\texttt{/fasilitas}.

\begin{figure}[H]
  \centering
  \includegraphics[width=0.9\textwidth]{../Img/Backend/screencapture-127-0-0-1-8000-admin-facilities-2026-06-16-17_54_00.png}
  \caption{Daftar Fasilitas di Panel Admin}
  \label{fig:facilities-index}
\end{figure}

\begin{figure}[H]
  \centering
  \includegraphics[width=0.9\textwidth]{../Img/Backend/screencapture-127-0-0-1-8000-admin-facilities-1-edit-2026-06-16-17_54_15.png}
  \caption{Form Edit Fasilitas}
  \label{fig:facilities-edit}
\end{figure}

\subsection{Panduan: Menambahkan Fasilitas}

\begin{enumerate}
  \item Buka menu \textbf{Fasilitas} dan klik tambah.
  \item Isi nama fasilitas (contoh: ``Laboratorium Komputer''), kapasitas, kondisi,
        dan deskripsi lengkap.
  \item Unggah foto fasilitas.
  \item Klik \textbf{Simpan}.
\end{enumerate}

% =====================================================================
%  BAB 4: MANAJEMEN MEDIA, GALERI, SLIDER & NAVIGASI
%  (Media Library, Downloads, Galleries, Sliders, Menus, Pages)
% =====================================================================

\chapter{Manajemen Media, Galeri, Slider \& Navigasi}

Modul ini mengatur aset visual website (gambar), dokumen publik yang bisa diunduh,
album foto, \textit{carousel} halaman depan, dan struktur navigasi website.

% -------------------------------------------------------------------
\section{Modul Pustaka Media (\textit{Media Library})}
% -------------------------------------------------------------------

Pustaka Media adalah repositori terpusat untuk semua berkas gambar yang diunggah
ke dalam sistem. Gambar yang diunggah secara otomatis dikonversi ke format
\textit{.webp} untuk mempercepat pemuatan halaman website.

\begin{figure}[H]
  \centering
  \includegraphics[width=0.9\textwidth]{../Img/Backend/screencapture-127-0-0-1-8000-admin-media-2026-06-16-17_52_59.png}
  \caption{Tampilan Pustaka Media --- Grid Gambar}
  \label{fig:media-index}
\end{figure}

\subsection{Cara Mengunggah Gambar}

\begin{enumerate}
  \item Klik menu \textbf{Media} di \textit{sidebar}.
  \item Klik tombol \textbf{Unggah Gambar} (atau seret file ke zona \textit{drop}).
  \item Pilih file dari komputer Anda. Format yang didukung: JPG, JPEG, PNG, GIF, WEBP.
  \item Gambar akan terproses dan muncul di kisi (\textit{grid}) galeri media.
\end{enumerate}

\subsection{Cara Menghapus Gambar}

\begin{enumerate}
  \item Arahkan kursor ke gambar yang ingin dihapus.
  \item Klik ikon \textbf{Hapus} (tempat sampah) yang muncul.
  \item Konfirmasi penghapusan.
\end{enumerate}

\begin{warnbox}[Proteksi Integritas Data]
  Sistem akan \textbf{menolak penghapusan} gambar yang masih tertaut ke konten aktif
  (misalnya gambar yang digunakan sebagai \textit{thumbnail} berita atau lampiran
  di modul Unduhan). Ini adalah fitur proteksi untuk mencegah \textit{broken links}
  di halaman publik. Hapus atau ganti keterkaitannya terlebih dahulu sebelum
  menghapus gambar dari pustaka media.
\end{warnbox}

% -------------------------------------------------------------------
\section{Modul Unduhan (\textit{Downloads})}
% -------------------------------------------------------------------

Modul Unduhan memungkinkan Admin menyediakan berkas dokumen (PDF, Word, Excel, dll.)
yang dapat diunduh oleh pengunjung publik melalui halaman \texttt{/unduhan}.

\begin{figure}[H]
  \centering
  \includegraphics[width=0.9\textwidth]{../Img/Backend/screencapture-127-0-0-1-8000-admin-downloads-2026-06-16-17_55_05.png}
  \caption{Daftar Berkas Unduhan di Panel Admin}
  \label{fig:downloads-index}
\end{figure}

\subsection{Skenario: Menyediakan Brosur PPDB}

Panitia PPDB meminta Admin untuk menyediakan brosur pendaftaran dalam format PDF
agar calon siswa bisa mengunduhnya dari website.

\subsection{Panduan: Menambahkan Berkas Unduhan}

\begin{enumerate}
  \item Buka menu \textbf{Unduhan} dan klik tombol tambah.
  \item Isi \textbf{Judul Dokumen} (contoh: ``Brosur PPDB Tahun Ajaran 2026/2027'').
  \item Isi \textbf{Deskripsi Singkat} (1-2 kalimat) tentang isi dokumen.
  \item Klik \textbf{Pilih File} pada kolom \textbf{File Dokumen (PDF, Word, Excel, dll.)}
        untuk mengunggah berkas dari komputer. Format yang diizinkan:
        \texttt{pdf, doc, docx, xls, xlsx, ppt, pptx, zip, rar}.
        Ukuran maksimal: \textbf{20 MB}.
  \item Centang opsi \textbf{Tersedia untuk Publik} agar berkas dapat diakses oleh
        pengunjung tanpa harus login.
        Jika tidak dicentang, berkas hanya tersimpan di sistem dan tidak tampil di halaman publik.
  \item Klik tombol hitam \textbf{Simpan}.
\end{enumerate}

\begin{notebox}[Mengganti File Lama]
  Saat mengedit unduhan yang sudah ada, Anda akan melihat informasi file aktif
  (nama file dan ukurannya dalam KB) beserta tautan \textbf{Lihat File} untuk verifikasi.
  Untuk mengganti file, cukup pilih file baru melalui tombol \textit{file picker}.
  Sistem akan menggantikan file lama secara otomatis.
\end{notebox}

\subsection{Penanganan Masalah: Unduhan}

\begin{itemize}
  \item \textbf{File gagal diunggah}: Periksa apakah ekstensi file termasuk dalam
        daftar yang diizinkan. File berekstensi \texttt{.php}, \texttt{.exe}, atau
        skrip berbahaya akan \textbf{ditolak secara otomatis} oleh sistem.
  \item \textbf{Unduhan tidak muncul di halaman publik}: Pastikan opsi
        \textbf{Tersedia untuk Publik} sudah dicentang.
\end{itemize}

% -------------------------------------------------------------------
\section{Modul Galeri (\textit{Galleries})}
% -------------------------------------------------------------------

Modul Galeri mengatur album foto kegiatan sekolah yang tampil di halaman \texttt{/galeri}.
Setiap album memiliki gambar \textit{cover} dan koleksi foto di dalamnya.

\begin{figure}[H]
  \centering
  \includegraphics[width=0.9\textwidth]{../Img/Backend/screencapture-127-0-0-1-8000-admin-galleries-2026-06-16-17_52_34.png}
  \caption{Daftar Album Galeri di Panel Admin}
  \label{fig:galleries-index}
\end{figure}

\begin{figure}[H]
  \centering
  \includegraphics[width=0.9\textwidth]{../Img/Backend/screencapture-127-0-0-1-8000-admin-galleries-2-edit-2026-06-16-17_52_48.png}
  \caption{Form Edit Album Galeri dan Manajemen Foto}
  \label{fig:galleries-edit}
\end{figure}

\subsection{Panduan: Membuat Album Galeri}

\begin{enumerate}
  \item Buka menu \textbf{Galeri} dan klik tambah.
  \item Isi \textbf{Judul} album (contoh: ``Hari Olahraga Nasional 2025'').
  \item Tentukan \textbf{Tanggal Publish} untuk album ini.
  \item Tulis \textbf{Deskripsi} singkat tentang album.
  \item Atur \textbf{Status}: \textit{Draft} atau \textit{Published}.
  \item Atur \textbf{Urutan} tampil di halaman galeri publik (angka kecil = tampil lebih awal).
  \item Pilih gambar \textbf{Cover Galeri} menggunakan komponen \textit{image-upload}.
  \item Klik tombol \textbf{Simpan}.
  \item Setelah album tersimpan, Anda bisa kembali ke halaman \textbf{Edit} album
        untuk menambahkan foto-foto individual ke dalam album tersebut
        melalui panel \textbf{Kelola Foto Album}.
\end{enumerate}

% -------------------------------------------------------------------
\section{Modul Slider (\textit{Sliders})}
% -------------------------------------------------------------------

Slider mengatur gambar-gambar yang berputar otomatis (\textit{carousel}) di bagian
\textit{Hero} halaman depan website. Setiap slide adalah sebuah foto \textit{landscape}
beresolusi tinggi.

\begin{figure}[H]
  \centering
  \includegraphics[width=0.9\textwidth]{../Img/Backend/screencapture-127-0-0-1-8000-admin-sliders-2026-06-16-17_53_13.png}
  \caption{Daftar Slide di Panel Admin}
  \label{fig:sliders-index}
\end{figure}

\subsection{Panduan: Menambahkan Slide Baru}

\begin{enumerate}
  \item Buka menu \textbf{Sliders} dan klik tambah.
  \item Tentukan \textbf{Urutan} slide (angka 1 = tampil pertama).
  \item Centang \textbf{Aktif} agar slide ini tampil di halaman depan.
        Slide yang tidak dicentang (\textit{inactive}) tidak akan ditampilkan
        meskipun tersimpan di sistem.
  \item Unggah \textbf{Gambar Slider} --- gunakan foto beresolusi tinggi
        dengan rasio \textbf{16:9} (minimal 1920×1080 px) agar tidak terpotong
        saat ditampilkan di berbagai ukuran layar.
  \item Klik tombol \textbf{Simpan}.
\end{enumerate}

\begin{notebox}[Zona Aman Gambar Slider]
  Letakkan elemen penting (teks, logo) pada bagian tengah gambar (\textit{safe zone})
  karena sistem akan memotong sisi gambar saat ditampilkan di layar ponsel.
\end{notebox}

% -------------------------------------------------------------------
\section{Modul Menu Navigasi (\textit{Menus})}
% -------------------------------------------------------------------

Menu Navigasi mengatur struktur tautan yang muncul di \textit{navbar} (header)
dan \textit{footer} website. Mendukung hierarki \textit{dropdown} (sub-menu) dan
pengaturan target tab.

\begin{figure}[H]
  \centering
  \includegraphics[width=0.9\textwidth]{../Img/Backend/screencapture-127-0-0-1-8000-admin-menus-1-edit-2026-06-16-17_53_28.png}
  \caption{Halaman Kelola Menu Navigasi dan Item-nya}
  \label{fig:menus-edit}
\end{figure}

\subsection{Panduan: Menambahkan Item Menu}

\begin{enumerate}
  \item Buka menu \textbf{Menus} lalu pilih menu yang ingin dikelola (misal: ``Menu Utama'').
  \item Pada panel \textbf{Tambah Item Menu}, isi:
        \begin{itemize}
          \item \textbf{Label} --- teks yang muncul di navigasi (contoh: ``Tentang Kami'')
          \item \textbf{URL} --- alamat tujuan (contoh: \texttt{/profil} untuk halaman internal,
                atau URL lengkap untuk tautan eksternal)
          \item \textbf{Target}:
                \begin{itemize}
                  \item \textit{Tab Sama} (\texttt{\_self}) --- buka di tab yang sama
                  \item \textit{Tab Baru} (\texttt{\_blank}) --- buka di tab baru
                        (disarankan untuk tautan eksternal)
                \end{itemize}
          \item \textbf{Parent} --- pilih item induk jika ini adalah sub-menu
                (misal: ``Profil'' sebagai \textit{parent} dari ``Visi Misi'')
          \item \textbf{Urutan} --- angka urut dalam grup yang sama
        \end{itemize}
  \item Klik tombol \textbf{Tambah Item}.
  \item Item yang baru ditambahkan akan muncul di daftar di bawahnya.
        Anda bisa mengubah urutan dengan fitur \textit{drag-and-drop} menggunakan
        ikon pegangan (\textit{drag handle}) di sisi kiri setiap baris.
\end{enumerate}

\subsection{Penanganan Masalah: Menu}

\begin{itemize}
  \item \textbf{Perubahan menu tidak terlihat di website}: Sistem mungkin menggunakan
        \textit{cache} untuk menu. Coba bersihkan \textit{cache} aplikasi melalui
        terminal dengan perintah \lstinline{php artisan cache:clear}, atau lakukan
        \textit{hard refresh} di peramban.
  \item \textbf{\textit{Dropdown} tidak muncul}: Pastikan item sub-menu sudah memiliki
        \textbf{Parent} yang benar. Tanpa \textit{parent}, item akan muncul sebagai
        menu utama, bukan sebagai \textit{dropdown}.
\end{itemize}

% =====================================================================
%  BAB 5: KONFIGURASI SISTEM & KUSTOMISASI LANJUTAN
%  (Settings, Users, Pods)
% =====================================================================

\chapter{Konfigurasi Sistem \& Kustomisasi Lanjutan}

Modul ini mencakup pengaturan global website, manajemen akun pengguna,
dan fitur \textit{Pods} (tipe konten kustom). Akses modul ini
\textbf{dibatasi hanya untuk Super Admin dan Admin}.

% -------------------------------------------------------------------
\section{Pengaturan Umum (\textit{General Settings})}
% -------------------------------------------------------------------

Pengaturan Umum menyimpan seluruh identitas sekolah yang ditampilkan di
berbagai halaman publik website.

\begin{figure}[H]
  \centering
  \includegraphics[width=0.9\textwidth]{../Img/Backend/screencapture-127-0-0-1-8000-admin-settings-2026-06-16-17_55_34.png}
  \caption{Tab Pengaturan Umum --- Identitas Sekolah}
  \label{fig:settings-general}
\end{figure}

\subsection{Panduan: Mengisi Identitas Sekolah}

\begin{enumerate}
  \item Klik menu \textbf{Pengaturan} (\textit{Settings}) di \textit{sidebar}.
        Anda langsung berada di tab \textbf{Umum}.
  \item Isi bagian \textbf{Identitas Sekolah}:
        \begin{itemize}
          \item \textbf{Nama Sekolah} --- nama resmi lengkap (wajib diisi)
          \item \textbf{Tagline} --- slogan atau motto sekolah
          \item \textbf{Deskripsi} --- paragraf perkenalan singkat
          \item \textbf{Alamat} --- alamat lengkap sekolah
          \item \textbf{Telepon} dan \textbf{Email}
          \item \textbf{Tahun Berdiri} dan \textbf{NPSN}
          \item \textbf{Akreditasi} dan \textbf{Nilai Akreditasi}
          \item \textbf{Status Sekolah} (Negeri/Swasta)
          \item \textbf{Luas Tanah (m2)}, \textbf{Luas Bangunan (m2)},
                \textbf{Waktu Belajar}, \textbf{Ruang Kelas}, \textbf{Program Studi}
        \end{itemize}
  \item Isi bagian \textbf{Logo \& Struktur}:
        \begin{itemize}
          \item \textbf{Logo Sekolah} --- digunakan juga di header, tab browser,
                dan favicon. Gunakan PNG berlatar transparan.
          \item \textbf{Gambar Struktur Organisasi} --- ditampilkan di halaman
                \texttt{/struktur-organisasi}.
        \end{itemize}
  \item Isi bagian \textbf{Profil Sekolah}:
        \begin{itemize}
          \item \textbf{Ringkasan Profil}, \textbf{Sejarah}, \textbf{Visi},
                \textbf{Misi} (baris baru per poin misi)
          \item \textbf{Video Profil (Link YouTube)} --- tempel URL penuh video,
                contoh: \texttt{https://www.youtube.com/watch?v=xxxxxxxx}.
                Video akan ditampilkan di beranda.
        \end{itemize}
  \item Isi bagian \textbf{Struktur \& Layanan}:
        \begin{itemize}
          \item \textbf{Deskripsi Struktur Organisasi}
          \item \textbf{Ringkasan Layanan} dan \textbf{Daftar Layanan}
                (baris baru per poin layanan)
        \end{itemize}
  \item Isi bagian \textbf{Statistik Beranda} --- angka ini ditampilkan di
        halaman depan sebagai \textit{stat counter}:
        \begin{itemize}
          \item \textbf{Total Guru \& Staf}
          \item \textbf{Siswa Aktif}
          \item \textbf{Program Unggulan}
          \item \textbf{Total Prestasi}
        \end{itemize}
  \item Isi bagian \textbf{Keamanan \& Akses}:
        \begin{itemize}
          \item \textbf{URL Login Admin Kustom} --- ganti \texttt{login} dengan
                kata rahasia Anda (misal: \texttt{portal-guru}).
                URL Login akan menjadi \texttt{/portal-guru}.
        \end{itemize}
  \item Klik tombol \textbf{Simpan Perubahan}.
\end{enumerate}

\begin{warnbox}[KRITIS: URL Login Kustom]
  Jika Anda mengubah \textbf{URL Login Admin Kustom} dan lupa kata rahasianya,
  Anda tidak akan bisa mengakses halaman login. Satu-satunya cara pemulihan adalah
  melalui \textit{database manager} (phpMyAdmin/Adminer) dengan mengubah nilai
  \texttt{admin\_login\_url} di tabel \texttt{settings} secara langsung.
  \textbf{Catat dan simpan URL rahasia ini di tempat yang aman!}
\end{warnbox}

% -------------------------------------------------------------------
\section{Pengaturan Kontak (\textit{Contact Settings})}
% -------------------------------------------------------------------

\begin{figure}[H]
  \centering
  \includegraphics[width=0.9\textwidth]{../Img/Backend/screencapture-127-0-0-1-8000-admin-settings-contact-2026-06-16-17_55_43.png}
  \caption{Tab Pengaturan Kontak --- WhatsApp dan Sosial Media}
  \label{fig:settings-contact}
\end{figure}

Navigasikan ke tab \textbf{Kontak} untuk mengisi:
\begin{itemize}
  \item \textbf{Maps Embed URL} --- URL \textit{iframe} Google Maps (untuk halaman \texttt{/kontak})
  \item \textbf{WhatsApp Number} --- nomor WhatsApp lengkap dengan kode negara (contoh: \texttt{628123456789})
  \item \textbf{Jam Operasional}
  \item Tautan media sosial: \textbf{Facebook URL}, \textbf{Instagram URL},
        \textbf{YouTube URL}, \textbf{Twitter/X URL}
\end{itemize}

Klik \textbf{Simpan Perubahan} setelah selesai.

% -------------------------------------------------------------------
\section{Pengaturan Tampilan (\textit{Appearance Settings})}
% -------------------------------------------------------------------

\begin{figure}[H]
  \centering
  \includegraphics[width=0.9\textwidth]{../Img/Backend/screencapture-127-0-0-1-8000-admin-settings-appearance-2026-06-16-17_55_58.png}
  \caption{Tab Pengaturan Tampilan --- Tema dan Tipografi}
  \label{fig:settings-appearance}
\end{figure}

Navigasikan ke tab \textbf{Tampilan} untuk mengatur:
\begin{itemize}
  \item \textbf{Warna Utama} dan \textbf{Warna Sekunder} --- kode warna HEX
        (contoh: \texttt{\#0ea5e9})
  \item \textbf{Jumlah Post per Halaman} --- berapa artikel yang ditampilkan
        per halaman di \texttt{/berita} (rentang: 1--60)
  \item \textbf{Font Utama (Teks Standar)} --- pilih dari daftar Google Fonts:
        \textit{Inter, Roboto, Poppins, Montserrat, Open Sans, Lato, Nunito,
        Outfit, Plus Jakarta Sans, Public Sans}
  \item \textbf{Font Judul (Aksen/Display)} --- pilih dari: \textit{Space Grotesk,
        Oswald, Bebas Neue, Raleway, DM Sans, Syne, Clash Display, Righteous, Cinzel}
  \item \textbf{Font Formal (Serif)} --- pilih dari: \textit{Playfair Display, Lora,
        Merriweather, PT Serif, Noto Serif, Crimson Text, EB Garamond}
\end{itemize}

Klik \textbf{Simpan Perubahan} setelah selesai.

% -------------------------------------------------------------------
\section{Pengaturan SEO}
% -------------------------------------------------------------------

\begin{figure}[H]
  \centering
  \includegraphics[width=0.9\textwidth]{../Img/Backend/screencapture-127-0-0-1-8000-admin-settings-seo-2026-06-16-17_56_11.png}
  \caption{Tab Pengaturan SEO Global}
  \label{fig:settings-seo}
\end{figure}

Navigasikan ke tab \textbf{SEO} untuk mengisi:
\begin{itemize}
  \item \textbf{Deskripsi Meta Default} --- deskripsi yang dipakai jika halaman
        tidak memiliki \textit{meta description} sendiri (maks. 320 karakter)
  \item \textbf{Google Analytics ID} --- kode pelacakan GA4
        (contoh: \texttt{G-XXXXXXXXXX})
  \item \textbf{Google Site Verification} --- kode verifikasi kepemilikan situs
        di Google Search Console
  \item \textbf{Robots.txt} --- isi aturan perayapan mesin pencari secara manual
\end{itemize}

% -------------------------------------------------------------------
\section{Pengaturan PPDB}
% -------------------------------------------------------------------

\begin{figure}[H]
  \centering
  \includegraphics[width=0.9\textwidth]{../Img/Backend/screencapture-127-0-0-1-8000-admin-settings-ppdb-2026-06-16-17_56_25.png}
  \caption{Tab Pengaturan PPDB --- Pendaftaran Peserta Didik Baru}
  \label{fig:settings-ppdb}
\end{figure}

Navigasikan ke tab \textbf{PPDB} untuk mengatur penerimaan siswa baru:
\begin{enumerate}
  \item Centang \textbf{Aktifkan Halaman PPDB} untuk memunculkan banner dan tombol
        ``Daftar PPDB'' di header website publik dan mengaktifkan halaman \texttt{/ppdb}.
  \item Isi \textbf{Link Pendaftaran (Opsional)} --- URL formulir pendaftaran eksternal
        (Google Form, website PPDB pusat, dll.).
  \item Isi \textbf{Tanggal Mulai (Opsional)} dan \textbf{Tanggal Berakhir (Opsional)}
        untuk menampilkan countdown atau informasi rentang waktu pendaftaran.
  \item Isi \textbf{Pengumuman / Informasi Tambahan PPDB} --- teks bebas untuk
        menampilkan syarat, narahubung, atau info gelombang pendaftaran.
  \item Klik tombol \textbf{Simpan Pengaturan}.
\end{enumerate}

% -------------------------------------------------------------------
\section{Manajemen Pengguna (\textit{Users})}
% -------------------------------------------------------------------

\begin{figure}[H]
  \centering
  \includegraphics[width=0.9\textwidth]{../Img/Backend/screencapture-127-0-0-1-8000-admin-users-2026-06-16-17_55_20.png}
  \caption{Daftar Pengguna Admin}
  \label{fig:users-index}
\end{figure}

\subsection{Panduan: Menambahkan Pengguna Baru}

\begin{enumerate}
  \item Klik menu \textbf{Pengguna} (\textit{Users}) di \textit{sidebar}.
  \item Klik tombol tambah.
  \item Isi \textbf{Nama} lengkap pengguna.
  \item Isi \textbf{Email} yang akan digunakan sebagai \textit{username} untuk login.
  \item Isi \textbf{Password} --- kata sandi sementara. Pengguna disarankan
        mengganti kata sandinya setelah login pertama kali.
        Saat mengedit pengguna yang sudah ada, \textbf{kosongkan kolom ini}
        jika tidak ingin mengubah kata sandi.
  \item Pilih \textbf{Role} dari \textit{dropdown}:
        \begin{itemize}
          \item \textit{Super Admin} --- akses penuh
          \item \textit{Admin} --- akses konten dan media
          \item \textit{Editor} --- akses tulis berita saja
        \end{itemize}
  \item Pastikan kotak \textbf{Aktif} dicentang agar pengguna dapat login.
  \item Klik tombol \textbf{Simpan}.
\end{enumerate}

\begin{notebox}[Menonaktifkan Pengguna]
  Untuk mencabut akses pengguna tanpa menghapus akunnya, \textbf{hapus centang}
  dari kotak \textbf{Aktif} dan simpan. Pengguna yang tidak aktif tidak dapat login
  ke panel admin.
\end{notebox}

% -------------------------------------------------------------------
\section{Tipe Konten Kustom (\textit{Pods})}
% -------------------------------------------------------------------

Fitur \textit{Pods} adalah mekanisme canggih yang memungkinkan Super Admin
membangun modul konten baru secara dinamis tanpa membutuhkan pemrograman.

\begin{figure}[H]
  \centering
  \includegraphics[width=0.9\textwidth]{../Img/Backend/screencapture-127-0-0-1-8000-admin-pods-2026-06-16-17_56_39.png}
  \caption{Halaman Manajemen Pods (Tipe Konten Kustom)}
  \label{fig:pods-index}
\end{figure}

\subsection{Skenario: Membuat Modul Testimoni Alumni}

Sekolah ingin menampilkan \textit{gallery} testimoni dari alumni terbaik di halaman depan.
Tidak ada modul bawaan untuk ini, maka Super Admin membuat \textit{Pod} baru.

\subsection{Panduan: Membuat Pod Baru}

\begin{enumerate}
  \item Buka menu \textbf{Pods} dan klik tambah.
  \item Isi \textbf{Nama Pod} (contoh: ``Testimoni Alumni'') dan deskripsinya.
  \item Simpan Pod. Sistem akan membuat struktur basis data secara otomatis.
  \item Klik \textbf{Kelola Field} pada Pod yang baru dibuat untuk mendefinisikan
        kolom-kolom datanya:
        \begin{itemize}
          \item Tambahkan field \textbf{Nama Lengkap} (tipe: Text)
          \item Tambahkan field \textbf{Angkatan} (tipe: Number)
          \item Tambahkan field \textbf{Kesan \& Pesan} (tipe: Textarea)
          \item Tambahkan field \textbf{Foto} (tipe: Image)
        \end{itemize}
  \item Setelah field tersimpan, buka menu \textbf{Pods $\rightarrow$ Testimoni Alumni $\rightarrow$ Entri}
        untuk mulai mengisi data testimoni alumni.
\end{enumerate}

\subsection{Menampilkan Data Pod di Konten Berita}

Data dari \textit{Pod} dapat disisipkan ke dalam konten berita menggunakan \textit{shortcode}
khusus di editor Quill:

\begin{lstlisting}
[pod:testimoni-alumni limit=5]
\end{lstlisting}

Sistem akan me-\textit{render} 5 entri testimoni terbaru dari Pod tersebut
secara dinamis di dalam konten artikel.

\subsection{Penanganan Masalah: Pods}

\begin{itemize}
  \item \textbf{\textit{Shortcode} tidak me-\textit{render}}:  Pastikan \textit{slug}
        Pod yang digunakan dalam \textit{shortcode} sama persis dengan \textit{slug}
        di sistem. Cek di halaman daftar Pods.
  \item \textbf{Data entri Pod tidak tersimpan}: Pastikan semua field yang
        didefinisikan sebagai \textit{required} sudah terisi sebelum menyimpan.
\end{itemize}


\end{document}
